% secciones/conclusiones.tex

\section{Conclusiones}

\subsection{Sobre el Data Warehouse retail con TPC-DS}

La implementación de un Data Warehouse basado en el benchmark \tpcds{} sobre Amazon EMR
permitió simular un entorno retail realista con millones de transacciones distribuidas
en cinco tablas relacionadas mediante un esquema estrella. La generación de 10 GB de datos
con \texttt{tpcds-kit} y su carga en HDFS demostró ser un proceso reproducible y
escalable, válido para evaluar tecnologías de procesamiento distribuido bajo condiciones
controladas.

\subsection{Sobre Apache Hive}

HiveQL demostró ser una herramienta madura y expresiva para consultas analíticas batch.
Su integración nativa con HDFS y el soporte de funciones de ventana como \texttt{RANK}
y \texttt{QUALIFY} permitieron implementar las nueve consultas del proyecto con código
SQL legible y mantenible. Su principal limitación es el tiempo de ejecución en consultas
complejas con múltiples joins, donde el motor Tez introduce etapas de materialización
en disco que incrementan la latencia.

\subsection{Sobre Apache Spark}

Apache Spark demostró un rendimiento superior en el benchmark comparativo, especialmente
en consultas que involucran múltiples joins y funciones de ventana sobre el dataset de
10 GB. El procesamiento en memoria y el optimizador Catalyst reducen significativamente
las operaciones de escritura en disco respecto a Hive. Adicionalmente, la API de PySpark
ofrece mayor flexibilidad para integrar lógica de negocio, como se evidenció en la
implementación de la capa agéntica, donde la \texttt{SparkSession} se reutiliza
entre consultas dinámicas generadas por el modelo LLM.

\subsection{Sobre el análisis agéntico}

La implementación de la capa agéntica basada en el paradigma de skills demostró que
es posible democratizar el acceso a un Data Warehouse sin requerir conocimiento de SQL
por parte del usuario final. El uso de la API de Gemini para generación automática de
consultas permitió responder preguntas en lenguaje natural con resultados correctos
en la mayoría de los casos evaluados.

Sin embargo, el enfoque agéntico presenta limitaciones inherentes: la calidad del SQL
generado depende de la claridad de la pregunta y de la completitud del esquema provisto
en el prompt. Preguntas ambiguas o que requieren joins no triviales pueden producir
resultados incorrectos sin que el sistema lo detecte automáticamente.

\subsection{Recomendaciones}

\begin{itemize}[noitemsep]
    \item Para cargas analíticas batch programadas, \hive{} es suficiente y tiene
          menor overhead de configuración.
    \item Para análisis exploratorio interactivo y pipelines dinámicos, \spark{} es
          la tecnología recomendada por su rendimiento y flexibilidad de API.
    \item El análisis agéntico es adecuado para usuarios no técnicos con preguntas
          directas sobre el negocio, pero debe complementarse con validación del SQL
          generado en entornos de producción.
    \item Para escalar el sistema agéntico, se recomienda agregar un paso de validación
          del SQL generado antes de ejecutarlo en Spark, y mantener un historial de
          preguntas y respuestas para mejorar el clasificador de intención.
\end{itemize}
